\chapter{Fenwick and segment trees}\label{ch:fenseg}
\index{Fenwick tree}\index{segment tree}\index{lazy propagation}

Prefix sums (Chapter~5) answer range queries once; this chapter
answers them \emph{under mutation}. Two structures cover the whole
family: the Fenwick tree for prefix-shaped questions, the segment
tree for everything else --- arbitrary monoids, range updates with
lazy tags, and order statistics by descent. After this chapter you
will pick the right structure off the update/query shape, write
each template, and know the two subtle failure modes that corrupt
silent answers.

\section{The range-query family}\label{sec:family}\index{range query}
Four shapes appear: point-update/range-query (the classic),
range-update/point-query (the difference trick), and
range-update/range-query (lazy). A decision table: prefix-ordered
aggregate and tight memory $\to$ Fenwick; arbitrary monoid or
descent queries $\to$ segment tree.

\section{The Fenwick tree}\label{sec:fenwick}\index{Fenwick tree!lowbit}
$t[i]$ stores the sum of the range $[i-\mathrm{lowbit}(i)+1,\ i]$ ---
a staircase of power-of-two lengths. Prefix sum walks down by
$i \mathrel{-}= \mathrm{lowbit}(i)$, update walks up by
$i \mathrel{+}= \mathrm{lowbit}(i)$: $O(\log n)$ each, $n$ cells.
\begin{example}[CSES \emph{Dynamic Range Sum Queries}]\label{ex:rsq}
Point updates, sum queries.
\par\smallskip
\textbf{Strategy.} The problem is exactly Fenwick's shape; we
show the segtree solution too for the trade-off table.
\par\smallskip
\begin{solution}
\lstinputlisting{code/dynamic_rsq.cpp}
Real output (CSES sample, segtree version):
\begin{lstlisting}[numbers=none]
14
2
11
2
\end{lstlisting}
\end{solution}
\end{example}
\begin{figure}[ht]\centering
\begin{tikzpicture}[font=\tiny,x=0.4cm,y=0.3cm]
\foreach \i/\h in {1/1,2/2,3/1,4/4,5/1,6/2,7/1,8/8}{
  \draw[slate,fill=accentsoft!40] (\i-1,0) rectangle (\i,\h);}
\node[font=\tiny] at (4,-1) {height = lowbit: $t[i]$ covers $i-\mathrm{lowbit}(i)+1..i$};
\end{tikzpicture}

\caption{The lowbit staircase. Notice every index's covered length
is its lowest set bit --- the whole structure is index arithmetic.}
\label{fig:fenwick}
\end{figure}

\section{The segment tree}\label{sec:segtree}\index{segment tree}
A $2n$ array: leaves at $n{+}i$, node $v$ combines children
$2v,2v{+}1$. Any associative monoid works --- sum, min, gcd ---
because a query only ever reads $O(\log n)$ whole nodes.
\begin{example}[CSES \emph{Dynamic Range Minimum Queries}]\label{ex:rmq}
Point updates, \emph{minimum} queries.
\par\smallskip
\textbf{Strategy.} Swap the monoid: same tree, \texttt{min} combine.
A Fenwick can't do this --- min has no inverse, so prefixes don't
subtract.
\par\smallskip
\begin{solution}
\lstinputlisting{code/dynamic_rmq.cpp}
Real output (CSES sample):
\begin{lstlisting}[numbers=none]
2
1
3
\end{lstlisting}
\end{solution}
\end{example}
\begin{figure}[ht]\centering
% fig62-segtree.tex -- Figure 12.x: segment tree over 8 leaves; node v has children 2v, 2v+1.
% fstyle.tex -- shared TikZ styles for the figures of this book (palette from lifeos.sty).
\tikzset{%
  cell/.style={draw=ink, minimum width=7mm, minimum height=6.5mm, inner sep=0pt, font=\footnotesize, fill=white},
  cellh/.style={cell, fill=accentsoft, draw=accent},
  cella/.style={cell, fill=amberbg, draw=amber},
  cellg/.style={cell, fill=goodbg, draw=good},
  cellx/.style={cell, fill=graybg, draw=slate, text=slate},
  vtx/.style={circle, draw=ink, minimum size=6.5mm, inner sep=0pt, font=\footnotesize, fill=white},
  vtxh/.style={vtx, draw=accent, fill=accentsoft},
  vtxa/.style={vtx, draw=amber, fill=amberbg},
  vtxg/.style={vtx, draw=good, fill=goodbg},
  vtxd/.style={vtx, draw=slate, fill=graybg, text=slate},
  lab/.style={font=\footnotesize, text=slate},
  labA/.style={font=\footnotesize, text=accent},
  labM/.style={font=\footnotesize, text=amber},
  idx/.style={font=\tiny, text=slate},
  arr/.style={-{stealth}, thick, ink},
  arrA/.style={-{stealth}, thick, accent},
  arrM/.style={-{stealth}, thick, amber},
  arrG/.style={-{stealth}, thick, good},
  edge/.style={thick, ink},
  edgeA/.style={thick, accent},
  edgeM/.style={thick, amber},
}%
% \cellrow{prefix}{x0}{y}{style}{v0,v1,...}: a row of cells named prefix0, prefix1, ... (pitch 0.75cm)
\providecommand{\cellrow}[5]{\foreach \v [count=\i from 0] in {#5} {\node[#4] (#1\i) at (#2+\i*0.75,#3) {\v};}}%
% \cellidx{prefix}{n}{start}: index labels under cells prefix0 .. prefix(n-1)
\providecommand{\cellidx}[3]{\foreach \i in {0,...,\numexpr#2-1\relax} {\node[idx, below=0pt of #1\i] {\the\numexpr\i+#3\relax};}}%
% \fitfig{<tikzpicture>}: shrink the picture to \linewidth only when it would overflow.
\ifdefined\fitbox\else\newsavebox{\fitbox}\fi%
\providecommand{\fitfig}[1]{\sbox{\fitbox}{#1}\ifdim\wd\fitbox>\linewidth\resizebox{\linewidth}{!}{\usebox{\fitbox}}\else\usebox{\fitbox}\fi}%

\begin{tikzpicture}[font=\footnotesize]
\tikzset{sn/.style={draw=slate, fill=accentsoft, rounded corners=1pt, minimum height=7mm, minimum width=1.3cm, inner sep=1pt, align=center, font=\scriptsize}}
\foreach \v/\lo/\hi/\x/\y in {1/1/8/5.25/3.0, 2/1/4/2.625/2.0, 3/5/8/7.875/2.0, 4/1/2/1.3125/1.0, 5/3/4/3.9375/1.0, 6/5/6/6.5625/1.0, 7/7/8/9.1875/1.0}
  \node[sn] (n\v) at (\x*1.0,\y*1.0) {{[\lo,\hi]}\\[-1pt]{\tiny node \v}};
\foreach \v/\lo/\x in {1/1/0.65, 2/2/1.97, 3/3/3.28, 4/4/4.6, 5/5/5.9, 6/6/7.2, 7/7/8.5, 8/8/9.8}
  \node[cell, minimum width=1.0cm, minimum height=6mm, font=\tiny] (l\v) at (\x*1.0+0.0,0) {$a_\lo$};
\draw[edge] (n1)--(n2) (n1)--(n3) (n2)--(n4) (n2)--(n5) (n3)--(n6) (n3)--(n7);
\draw[edge] (n4)--(l1) (n4)--(l2) (n5)--(l3) (n5)--(l4) (n6)--(l5) (n6)--(l6) (n7)--(l7) (n7)--(l8);
\node[lab, anchor=west] at (0,-0.9) {array of size $2n$: node $v$ stores the combine of its segment; children are $2v$ and $2v{+}1$};
\end{tikzpicture}

\caption{Node $\leftrightarrow$ segment. Notice the query reads
\emph{nodes}, not elements --- each node is a whole covered range.}
\label{fig:segtree}
\end{figure}
\begin{figure}[ht]\centering
\begin{tikzpicture}[font=\tiny]
\draw[slate!50] (0,0) rectangle (4,0.6);
\node[font=\tiny] at (2,-0.35) {array of 8};
\node[draw=ink,fill=accent!40] at (1,0.3) {[1,4]};
\node[draw=ink,fill=accentsoft!40] at (2.5,0.3) {[5]};
\node[draw=ink,fill=accentsoft!40] at (3.5,0.3) {[6,8]};
\node[font=\tiny] at (2,-1) {query decomposes into $\le 2\log n$ nodes};
\end{tikzpicture}

\caption{Range decomposition. Notice a query peels at most two
nodes per level --- $\le 2\log n$ nodes total.}
\label{fig:decomp}
\end{figure}

\section{Lazy propagation}\label{sec:lazy}\index{lazy propagation}
A range update stops at covered nodes, tagging them; the tag
travels down only when a later query needs through. The canonical
pair: add-tag with sum-query.
\begin{example}[CSES \emph{Range Update Queries}]\label{ex:ruq}
Range updates, point queries.
\par\smallskip
\textbf{Strategy.} The \emph{difference trick} turns range-add into
two point-adds on a prefix array --- a Fenwick of differences. The
lazy segtree generalizes to range-range.
\par\smallskip
\begin{solution}
\lstinputlisting{code/range_update_queries.cpp}
Real output (CSES sample):
\begin{lstlisting}[numbers=none]
5
\end{lstlisting}
\end{solution}
\end{example}
\begin{figure}[ht]\centering
% fig64-lazy.tex -- Figure 12.x: lazy propagation; add +5 on [1,4] tags one node; a later query through it pushes the tag down.
% fstyle.tex -- shared TikZ styles for the figures of this book (palette from lifeos.sty).
\tikzset{%
  cell/.style={draw=ink, minimum width=7mm, minimum height=6.5mm, inner sep=0pt, font=\footnotesize, fill=white},
  cellh/.style={cell, fill=accentsoft, draw=accent},
  cella/.style={cell, fill=amberbg, draw=amber},
  cellg/.style={cell, fill=goodbg, draw=good},
  cellx/.style={cell, fill=graybg, draw=slate, text=slate},
  vtx/.style={circle, draw=ink, minimum size=6.5mm, inner sep=0pt, font=\footnotesize, fill=white},
  vtxh/.style={vtx, draw=accent, fill=accentsoft},
  vtxa/.style={vtx, draw=amber, fill=amberbg},
  vtxg/.style={vtx, draw=good, fill=goodbg},
  vtxd/.style={vtx, draw=slate, fill=graybg, text=slate},
  lab/.style={font=\footnotesize, text=slate},
  labA/.style={font=\footnotesize, text=accent},
  labM/.style={font=\footnotesize, text=amber},
  idx/.style={font=\tiny, text=slate},
  arr/.style={-{stealth}, thick, ink},
  arrA/.style={-{stealth}, thick, accent},
  arrM/.style={-{stealth}, thick, amber},
  arrG/.style={-{stealth}, thick, good},
  edge/.style={thick, ink},
  edgeA/.style={thick, accent},
  edgeM/.style={thick, amber},
}%
% \cellrow{prefix}{x0}{y}{style}{v0,v1,...}: a row of cells named prefix0, prefix1, ... (pitch 0.75cm)
\providecommand{\cellrow}[5]{\foreach \v [count=\i from 0] in {#5} {\node[#4] (#1\i) at (#2+\i*0.75,#3) {\v};}}%
% \cellidx{prefix}{n}{start}: index labels under cells prefix0 .. prefix(n-1)
\providecommand{\cellidx}[3]{\foreach \i in {0,...,\numexpr#2-1\relax} {\node[idx, below=0pt of #1\i] {\the\numexpr\i+#3\relax};}}%
% \fitfig{<tikzpicture>}: shrink the picture to \linewidth only when it would overflow.
\ifdefined\fitbox\else\newsavebox{\fitbox}\fi%
\providecommand{\fitfig}[1]{\sbox{\fitbox}{#1}\ifdim\wd\fitbox>\linewidth\resizebox{\linewidth}{!}{\usebox{\fitbox}}\else\usebox{\fitbox}\fi}%

\begin{tikzpicture}[font=\footnotesize]
\tikzset{sn/.style={draw=slate, fill=white, rounded corners=1pt, minimum height=7mm, minimum width=1.35cm, inner sep=1pt, font=\scriptsize, align=center},
         tg/.style={sn, draw=amber, fill=amber!25, very thick}}
% panel A: after update
\node[labA] at (2.0,3.1) {1. update $[1,4]\mathrel{+}=5$};
\node[sn] (r) at (2.0,2.3) {$[1,8]$\\sum $+20$};
\node[tg] (a) at (0.5,1.2) {$[1,4]$\\tag $+5$};
\node[sn] (b) at (3.5,1.2) {$[5,8]$};
\node[sn] (a1) at (-0.4,0) {$[1,2]$\\clean}; \node[sn] (a2) at (1.4,0) {$[3,4]$\\clean};
\draw[slate] (r)--(a) (r)--(b) (a)--(a1) (a)--(a2);
\node[labM, font=\scriptsize, align=center] at (1.5,-0.9) {the tag stays on $[1,4]$:\\ children are not touched};
% arrow
\draw[arrA, ultra thick] (4.9,1.2) -- (5.9,1.2) node[midway, above, labA] {query};
% panel B: after push
\node[labA] at (9.0,3.1) {2. a query descends through $[1,4]$};
\node[sn] (r2) at (9.0,2.3) {$[1,8]$};
\node[sn] (a_) at (7.5,1.2) {$[1,4]$\\tag cleared}; \node[sn] (b_) at (10.5,1.2) {$[5,8]$};
\node[tg] (a1_) at (6.6,0) {$[1,2]$\\tag $+5$}; \node[tg] (a2_) at (8.4,0) {$[3,4]$\\tag $+5$};
\draw[slate] (r2)--(a_) (r2)--(b_) (a_)--(a1_) (a_)--(a2_);
\node[labM, font=\scriptsize, align=left, anchor=west] at (9.4,0.0) {push: the tag moves one\\ level down only when needed};
\end{tikzpicture}

\caption{Tag deferral. Notice the invariant: a node's aggregate is
always correct; its descendants are behind by the tag.}
\label{fig:lazy}
\end{figure}

\section{Walking the tree}\label{sec:walk}\index{order statistic}
Descent queries need no new code: at node $v$, if the left child's
value decides, go left, else right. ``First room with enough
capacity'' and ``the $k$-th remaining element'' are the same walk.
\begin{example}[CSES \emph{Hotel Queries}]\label{ex:hotel}
Each group wants the first room holding them.
\par\smallskip
\textbf{Strategy.} Store subtree max; descend left whenever
left-max $\ge r$. One trap: the walk needs leaves on a
power-of-two grid --- pad the array up.
\par\smallskip
\begin{solution}
\lstinputlisting{code/hotel_queries.cpp}
Real output (CSES sample):
\begin{lstlisting}[numbers=none]
3 8 1 0 0 
\end{lstlisting}
\end{solution}
\end{example}
\begin{figure}[ht]\centering
\begin{tikzpicture}[font=\tiny]
\node[draw=slate] (r) at (0,0) {max 9};
\node[draw=slate,fill=accentsoft!40] (a) at (-1.6,-0.9) {max 7};
\node[draw=slate] (b) at (1.6,-0.9) {max 9};
\node[draw=slate,fill=accentsoft!40] (a1) at (-2.4,-1.8) {max 6};
\node[draw=slate] (a2) at (-0.9,-1.8) {max 7};
\draw[-{stealth},amber] (r)--(a);
\draw[-{stealth},amber] (a)--(a1);
\draw[slate] (r)--(b) (a)--(a2);
\node[font=\tiny] at (0,-2.6) {descend left while left-max $\ge r$};
\end{tikzpicture}

\caption{The descent. Notice the search is guided by the same
aggregate the query would read --- no separate index needed.}
\label{fig:walk}
\end{figure}

\section{Order statistics}\label{sec:orderstat}\index{k-th element}
The Fenwick of \emph{counts} supports $k$-th-element descent and
erases --- ``which position is the $k$-th survivor''.
\begin{example}[CSES \emph{List Removals}]\label{ex:removals}
Repeatedly remove the $k$-th remaining element.
\par\smallskip
\textbf{Strategy.} Fenwick over 1s; \texttt{lower\_bound(k)} finds
the position whose count-prefix hits $k$; erase with
\texttt{add(p,-1)}.
\par\smallskip
\begin{solution}
\lstinputlisting{code/list_removals.cpp}
Real output (CSES sample):
\begin{lstlisting}[numbers=none]
1 2 2 6 4
\end{lstlisting}
\end{solution}
\end{example}

\begin{exercise}[easy]\label{exr:four}
Which of the four update/query shapes is \emph{not} solvable by
a plain Fenwick + difference array?
\end{exercise}
\begin{exercise}[medium]\label{exr:minbit}
Explain exactly why a Fenwick can't answer range-min ---
what property of sum is being used?
\end{exercise}
\begin{exercise}[medium]\label{exr:pushorder}
In lazy propagation, why must the tag be pushed \emph{before}
reading the children's aggregates?
\end{exercise}
\begin{exercise}[medium]\label{exr:descinv}
State the loop invariant of the tree descent and use it to
argue the walk ends at the leftmost qualifying leaf.
\end{exercise}

\section{Chapter summary}
\begin{table}[ht]\centering
{\small
\begin{tabularx}{\linewidth}{@{}lYYY@{}}
\toprule
\tablehead\thead{Shape} & \thead{Structure} & \thead{Cost} & \thead{Limit}\\
\midrule
point-upd, prefix-q & Fenwick & $\log n$, $n$ cells & invertible only\\
point-upd, range-q & segtree & $\log n$, $2n$ & any monoid\\
range-upd, point-q & diff + Fenwick & $\log n$ & no range query\\
range-upd, range-q & lazy segtree & $\log n$, $4n$ & tag design\\
descent queries & walk on tree & $\log n$ & pow-2 grid\\
\bottomrule
\end{tabularx}}
\caption{Pick by the update/query pair, then by the monoid.}
\label{tab:fenseg}
\end{table}

\section{Common mistakes}
\index{common mistakes}
\begin{itemize}
\item \textbf{Fenwick on non-invertible ops.} Min/max have no
inverse --- prefix differences don't exist; use a segtree.
\item \textbf{Descent on a non-power-2 layout.} Iterative segtrees
interleave leaf ranges when $n$ isn't a power of two; the walk
silently lands on the wrong leaf --- pad to a power of two.
\item \textbf{Forgetting the push.} Reading a child under an
unpushed tag returns stale data --- always push on the way down.
\item \textbf{Off-by-one leaves.} $t[n+i]$ indexes position $i$,
not $i{+}1$ --- mixing conventions shifts every update.
\item \textbf{Tag/aggregate mismatch.} The lazy tag must compose
with \emph{itself} in the update's order; add-then-assign is a
different monoid than assign-then-add.
\end{itemize}

\section{Additional exercises}
\begin{exercise}[easy] CSES \emph{Static Range Sum Queries}. \end{exercise}
\begin{exercise}[easy] CSES \emph{Static Range Minimum Queries}. \end{exercise}
\begin{exercise}[easy] CSES \emph{Range Xor Queries}. \end{exercise}
\begin{exercise}[easy] USACO-G \emph{Haircut} --- inversions on BIT. \end{exercise}
\begin{exercise}[easy] USACO-G \emph{Balanced Photo}. \end{exercise}
\begin{exercise}[medium] CSES \emph{Dynamic Range Sum Queries}. \end{exercise}
\begin{exercise}[medium] CSES \emph{Dynamic Range Minimum Queries}. \end{exercise}
\begin{exercise}[medium] CSES \emph{Salary Queries} --- compress + descent. \end{exercise}
\begin{exercise}[medium] USACO-G \emph{Sleepy Cow Sorting} --- BIT of counts. \end{exercise}
\begin{exercise}[medium] USACO-G \emph{Why Did the Cow Cross the Road III}. \end{exercise}
\begin{exercise}[medium] Kattis \emph{Mega Inversions}. \end{exercise}
\begin{exercise}[hard] CSES \emph{Prefix Sum Queries}. \end{exercise}
\begin{exercise}[hard] CSES \emph{Range Updates and Sums} --- two-tag lazy. \end{exercise}
\begin{exercise}[hard] CSES \emph{Distinct Values Queries}. \end{exercise}
\begin{exercise}[hard] CF \emph{Irrigation}. \end{exercise}

\chapternotes
Fenwick's lowbit staircase and the iterative segtree layout follow
Laaksonen~\cite{laaksonen2024} and Sannemo~\cite{sannemo2025}; the
Al.Cash iterative segtree and descent patterns are the
cp-algorithms treatment~\cite{cpalgorithms} (with the
USACO-Guide lazy-propagation walkthrough~\cite{usacoguide}); Halim's
structure-selection advice shaped the decision
table~\cite{halimcp4}; archive problems are official
CSES/USACO~\cite{cses}.

\section*{Further reading}
\begin{itemize}
\item USACO Guide, ``Point Update Range Sum'' (intermediate) --- the Fenwick/segtree
point-update family of \S\ref{sec:family} with Gold problems. \newline\url{https://usaco.guide/gold/PURS}
\item cp-algorithms, ``Fenwick Tree'' (reference) ---
\S\ref{sec:fenwick} plus range-update tricks and the
\texttt{lower\_bound} descent of \S\ref{sec:orderstat}. \newline\url{https://cp-algorithms.com/data_structures/fenwick.html}
\item cp-algorithms, ``Segment Tree'' (reference) ---
\S\ref{sec:segtree}--\S\ref{sec:lazy} including lazy propagation in
full generality. \newline\url{https://cp-algorithms.com/data_structures/segment_tree.html}
\item USACO Guide, ``Segment Tree Beats'' (advanced) --- what lies beyond the
monoid rule of \S\ref{sec:segtree}: chmin-style range bounds. \newline\url{https://usaco.guide/adv/segtree-beats}
\item CSES 1648 \emph{Dynamic Range Sum Queries} and 1649
\emph{Dynamic Range Minimum Queries} (practice) ---
\ref{ex:rsq}/\ref{ex:rmq} verbatim. \newline\url{https://cses.fi/problemset/task/1648}
\item CSES 1143 \emph{Hotel Queries} and 1749 \emph{List Removals}
(practice) --- the
descent-on-tree problems of \S\ref{sec:walk}--\S\ref{sec:orderstat}. \newline\url{https://cses.fi/problemset/task/1143}\end{itemize}

\begin{rulebox}[Watch instead]
\smallskip\noindent\textbf{Best first.}
``Fenwick Tree (Binary Index Tree) --- Quick Tutorial and Source Code
Explanation'' --- Stable Sort, 9 min ---
\url{https://youtu.be/uSFzHCZ4E-8}. The lowbit trick
of \S\ref{sec:fenwick} line by line; uses \texttt{bit} array name for
the book's \texttt{t}, same indexing.
\smallskip\noindent
``Segment Tree (Implementation)'' --- Errichto Hard Algorithms,
79 min --- \url{https://youtu.be/2FShdqn-Oz8}.
Recursive segtree and lazy propagation coded live, covering
\S\ref{sec:segtree}--\S\ref{sec:lazy}; deeper and longer than the
chapter --- pairs with \texttt{lazyseg.h}. Access date 2026-10-10.
\end{rulebox}
